\section{Curvature, concentration, and exact susceptibilities}
\label{ac:section}

The attractive determinant measure has a geometric positivity property
before any mean-field limit is taken.  Its density is the canonical
volume of a positively curved metric on the anticanonical line bundle
of the full configuration space.  The resulting estimate controls
empirical observables uniformly in the particle number, including
shrinking projective-moment tubes.

Write $X=\mathbb P^q$, $D=q+1$, $L=-K_X=\mathcal O(D)$,
$N=N_k=\binom{Dk+q}{q}$, and let $\sigma$ be Fubini--Study
probability.  Give $L$ its invariant metric $h_0=e^{-\phi_0}$, and
choose an orthonormal basis of $H^0(X,L^k)$ for $(h_0^k,\sigma)$.
Its evaluation determinant is a nonzero section $\Delta_k$ of
$\boxtimes_{i=1}^N L^k$ on $M=X^N$.  Put
\[
 w_{k,\gamma}=\|\Delta_k\|_{h_0^{k\boxtimes N}}^{-2\gamma/k},
 \qquad
 Z_{k,\gamma}=\int_M w_{k,\gamma}\,d\sigma^{\otimes N},
 \qquad
 dP_{k,\gamma}=Z_{k,\gamma}^{-1}w_{k,\gamma}
                                      \,d\sigma^{\otimes N}.
\]
The value on the zero divisor is immaterial whenever the partition
function is finite.  Additive constants in $\phi_0$ or a constant
rescaling of $\Delta_k$ do not change $P_{k,\gamma}$.

The metric convention in the estimates below is explicit.  Set
\[
 G_X=i\partial\bar\partial\phi_0=2\pi\omega_0,
 \qquad \omega_0=D\omega_{\rm FS},
 \qquad G_M=\sum_{i=1}^N\pi_i^*G_X.
\]
Here $dd^c=i\partial\bar\partial/(2\pi)$.  All covector norms
and Hessian comparisons in this section use the Hermitian metric with
K\"ahler form $G_X$ or $G_M$, rather than $\omega_0$.

\subsection{The anticanonical variance estimate}

\begin{lemma}[Canonical-volume Poincar\'e inequality]
\label{ac:canonical-poincare}
Let $(M,G)$ be a connected compact K\"ahler manifold and let
$e^{-\Phi}$ be a smooth metric on $-K_M$ satisfying
$i\partial\bar\partial\Phi\ge aG$, where $a>0$.
Let $\mu_\Phi$ be its canonical volume, normalized to mass one.
Then every smooth complex-valued function $F$ satisfies
\begin{equation}
 \int_M|F-\mu_\Phi(F)|^2\,d\mu_\Phi
 \le \frac1a\int_M|\bar\partial F|_G^2\,d\mu_\Phi.
 \label{ac:canonical-poincare-eq}
\end{equation}
\end{lemma}

\begin{proof}
Regard $F$ as a section of
$K_M\otimes(-K_M)\simeq\mathcal O_M$.  In a holomorphic coordinate
chart, its squared $L^2$ norm as a $(-K_M)$-valued top-degree form is
$|F|^2e^{-\Phi}$ times coordinate volume: the metric factors from
$K_M$ and the auxiliary volume of $G$ cancel.  The analogous norm
of $\bar\partial F$ is
$|\bar\partial F|_G^2e^{-\Phi}$ times coordinate volume.

The Bochner--Kodaira identity on $(-K_M)$-valued $(\dim M,1)$-forms
has curvature term $[i\Theta(-K_M),\Lambda_G]\ge aI$.
Consequently the $L^2$ estimate for the minimal solution of
$\bar\partial U=\bar\partial F$ gives
\[
 \|U\|_{L^2(e^{-\Phi})}^2
 \le a^{-1}\|\bar\partial F\|_{L^2(G,e^{-\Phi})}^2.
\]
This is the compact K\"ahler form of H\"ormander's estimate
\cite{Hormander1965}; it also follows directly from the preceding
curvature lower bound and the Hilbert-space minimal-solution argument.
The holomorphic sections of $K_M\otimes(-K_M)$ are exactly the
constant functions.  Thus the minimal solution is
$U=F-\mu_\Phi(F)$.  Dividing both sides by the mass of the canonical
volume proves the assertion.
\end{proof}

\begin{proposition}[Attractive covariance and exponential tilts]
\label{ac:variance}
Suppose $0\le\gamma<1$ and $Z_{k,\gamma}<\infty$.
Then, for every smooth function $F$ on $X^N$,
\begin{equation}
 \operatorname{Var}_{P_{k,\gamma}}(F)
 \le\frac1{1-\gamma}
       \int|\bar\partial F|_{G_M}^2\,dP_{k,\gamma}.
 \label{ac:variance-eq}
\end{equation}
For a real $f\in C^\infty(X)$, put
\[
 A_f=\sup_X|\bar\partial f|_{G_X}^2,
 \qquad
 -C_fG_X\le i\partial\bar\partial f\le C_fG_X.
\]
Here $C_f$ is any nonnegative constant satisfying the displayed
inequalities.  If
$dP_{k,\gamma,t,f}\propto e^{t\sum_i f(x_i)}dP_{k,\gamma}$ and
$|t|C_f<1-\gamma$, then
\begin{equation}
 \operatorname{Var}_{P_{k,\gamma,t,f}}
          \left(\sum_{i=1}^N f(x_i)\right)
 \le\frac{NA_f}{1-\gamma-|t|C_f}.
 \label{ac:tilted-variance}
\end{equation}
Every one-particle marginal of $P_{k,\gamma}$ is exactly $\sigma$.
\end{proposition}

\begin{proof}
Write $\Phi_0=\sum_i\phi_0(x_i)$ in a local frame of $-K_M$ and
write $s$ for the corresponding holomorphic representative of
$\Delta_k$.  The unnormalized density is the canonical volume of
the singular metric
\[
 \Phi_\gamma=(1-\gamma)\Phi_0+\frac\gamma k\log|s|^2.
\]
For $\varepsilon>0$ use the smooth metric
\begin{equation}
 \Phi_{\gamma,\varepsilon}
  =(1-\gamma)\Phi_0+\frac\gamma k
                 \log\bigl(|s|^2+\varepsilon e^{k\Phi_0}\bigr).
 \label{ac:regularization}
\end{equation}
Under a change of line-bundle frame the expression in parentheses
changes by the appropriate squared transition factor, so this formula
defines a global metric.  The logarithm is plurisubharmonic: it is
the logarithmic sum of exponentials of the plurisubharmonic functions
$\log|s|^2$ and $\log\varepsilon+k\Phi_0$.  It is smooth even at
$s=0$.  Hence
\[
 i\partial\bar\partial\Phi_{\gamma,\varepsilon}
                          \ge(1-\gamma)G_M.
\]
Its canonical density, relative to $\sigma^{\otimes N}$ up to an
irrelevant constant, is
$(\|\Delta_k\|^2+\varepsilon)^{-\gamma/k}$.
As $\varepsilon\downarrow0$, these densities increase to
$w_{k,\gamma}$, and their integrals increase to $Z_{k,\gamma}<\infty$.
The resulting probability measures therefore converge in total
variation.  Apply Lemma~\ref{ac:canonical-poincare} to the regularized
metrics and pass to the limit.  Both the integrands on the right and
the first two moments on the left are bounded for a fixed smooth $F$.
This proves \eqref{ac:variance-eq} without applying a smooth theorem
directly to a singular metric.

Multiplication of the density by $e^{t\sum_i f(x_i)}$ changes the
metric in \eqref{ac:regularization} to
$\Phi_{\gamma,\varepsilon}-t\sum_i f(x_i)$.  Its curvature is at
least $(1-\gamma-|t|C_f)G_M$.  In the product metric,
\[
 \left|\bar\partial\sum_i f(x_i)\right|_{G_M}^2
             =\sum_i|\bar\partial f(x_i)|_{G_X}^2\le NA_f.
\]
The same approximation proves \eqref{ac:tilted-variance}.

Finally, the diagonal action of $SU(D)$ preserves the product reference
measure and the squared determinant norm.  Indeed, it acts unitarily
on $H^0(X,L^k)$, so the evaluation determinant changes only by the
determinant of a unitary basis-change matrix.  Its modulus is unchanged.
Every one-particle marginal is therefore $SU(D)$-invariant, and the
unique invariant probability on $X$ is $\sigma$.
\end{proof}

\subsection{An exact projective Ward identity}

\begin{proposition}[Sharp covariance of the first moment]
\label{ac:ward}
Suppose $\gamma<1$ and $Z_{k,\gamma}<\infty$, using the same
density formula also for negative $\gamma$.
For traceless Hermitian $D\times D$ matrices $H,K$, let
$f_H(z)=\operatorname{tr}(HP_z)$ and $S_H=\sum_i f_H(z_i)$.
Then
\begin{equation}
 \operatorname{Cov}_{P_{k,\gamma}}(S_H,S_K)
 =\frac{N\operatorname{tr}(HK)}{D(D+1)(1-\gamma)}.
 \label{ac:ward-covariance}
\end{equation}
In the attractive range $0\le\gamma<1$, the constant
$1/(1-\gamma)$ in \eqref{ac:variance-eq} is therefore optimal.
For two distinct labels,
\begin{equation}
 \mathbb E_{P_{k,\gamma}}\operatorname{tr}(P_{z_1}P_{z_2})
 =\frac1D+
   \frac{(D-1)\gamma}{D(N-1)(1-\gamma)}.
 \label{ac:pair-overlap}
\end{equation}
Equivalently, with $M_1=N^{-1}\sum_iP_{z_i}$,
\begin{equation}
 \mathbb E_{P_{k,\gamma}}
      \left\|M_1-\frac{I}{D}\right\|_{\rm HS}^2
       =\frac{D-1}{DN(1-\gamma)}.
 \label{ac:first-moment-variance}
\end{equation}
Consequently finiteness of the partition function in this range
requires $\gamma<1-1/N$.
\end{proposition}

\begin{proof}
For $A\in SL(D,\mathbb C)$ the projective map
$g_A(z)=Az/\|Az\|$ has Jacobian $\|Az\|^{-2D}$ relative to
$\sigma$.  The determinant of the $SL(D,\mathbb C)$ representation
on $\operatorname{Sym}^{Dk}\mathbb C^D$ is one.  Multihomogeneity
of the evaluation determinant therefore gives
\[
 \|\Delta_k(g_Az_1,\ldots,g_Az_N)\|^2
 =\|\Delta_k(z_1,\ldots,z_N)\|^2
                  \prod_i\|Az_i\|^{-2Dk}.
\]
Change variables by $g_A$ in the finite partition function.  Taking
$A=e^{tH/2}$ gives the exact identity
\begin{equation}
 \mathbb E_{P_{k,\gamma}}
 \exp\!\left[-D(1-\gamma)
       \sum_i\log(z_i^*e^{tH}z_i)\right]=1.
 \label{ac:ward-identity}
\end{equation}
For $t$ in a compact interval the exponent and its first two
derivatives are bounded, so differentiation is justified by the
already assumed finiteness of $Z_{k,\gamma}$.
The first derivative vanishes, and the second gives
\[
 D(1-\gamma)\mathbb E S_H^2
 =\mathbb E\sum_i
       \left[\operatorname{tr}(H^2P_{z_i})-f_H(z_i)^2\right].
\]
Every one-particle marginal is Haar.  Its second-moment identities are
\[
 \mathbb E_\sigma\operatorname{tr}(H^2P)=\frac{\operatorname{tr}H^2}{D},
 \qquad
 \mathbb E_\sigma f_H^2=\frac{\operatorname{tr}H^2}{D(D+1)}.
\]
Substitution and polarization prove \eqref{ac:ward-covariance}.
Direct differentiation of $f_H$ in the invariant metric gives
\[
 |\bar\partial f_H|_{G_X}^2
 =\frac1D\left[\operatorname{tr}(H^2P)-f_H^2\right].
\]
In the attractive range, equality therefore holds in
\eqref{ac:variance-eq} for $F=S_H$ whenever $H\ne0$, proving
sharpness.  The Ward-identity calculation itself is unchanged for
negative $\gamma$, where the partition function is finite because
the density is bounded on the compact configuration space.

Permutation symmetry and the known one-particle covariance imply
\[
 \mathbb E f_H(z_1)f_K(z_2)
 =\frac{\gamma\operatorname{tr}(HK)}
 {(N-1)(1-\gamma)D(D+1)}.
\]
Sum over a Hilbert--Schmidt orthonormal basis of traceless Hermitian
matrices and use
$\operatorname{tr}(PQ)=D^{-1}+\sum_a
\operatorname{tr}(H_aP)\operatorname{tr}(H_aQ)$ to obtain
\eqref{ac:pair-overlap}.  Summing \eqref{ac:ward-covariance} in the
same basis and dividing by $N^2$ proves
\eqref{ac:first-moment-variance}.  For two distinct coordinate labels
under a measure absolutely continuous with respect to
$\sigma^{\otimes N}$, one has
$\operatorname{tr}(P_{z_1}P_{z_2})<1$ almost surely.
The right side of \eqref{ac:pair-overlap} must therefore be strictly
less than one, which is equivalent to $\gamma<1-1/N$.
\end{proof}

\subsection{An unconditional attractive range}

\begin{lemma}[Integrability by monomial degeneration]
\label{ac:integrability}
For every $k\ge1$ the partition function is finite whenever
$0\le\gamma<1/D$.  More generally, every nonzero section of
$\mathcal O(m,\ldots,m)$ on $(\mathbb P^q)^N$ has complex
singularity exponent at least $1/m$.
\end{lemma}

\begin{proof}
Expand a nonzero section in the multihomogeneous monomial basis.
Choose integer weights on the homogeneous coordinates, independently
in each factor, such that one monomial in its finite support uniquely
minimizes the total weight.  Pullback by the corresponding toric
one-parameter subgroup, followed by multiplication by the scalar
power of the parameter that sets this minimum weight to zero, gives
a holomorphic family whose limit is that single monomial.  For a
nonzero parameter the global complex singularity exponent is unchanged:
the transformation is an automorphism and the scalar is nonzero.

The divisor of the limiting monomial is a sum of the coordinate
hyperplanes, each with multiplicity at most $m$.  These hyperplanes
form a simple normal-crossing divisor: in each projective factor at
most $q$ of its $q+1$ coordinate hyperplanes meet.  In local coordinates
integrability of a monomial is the product of the conditions
$c\alpha_j<1$; its global exponent is therefore at least $1/m$.
Lower semicontinuity of complex singularity exponents
\cite{geo_demailly_kollar2001} shows that the original section has
exponent at least that of this limit.  Smooth positive metric and
volume factors do not affect local integrability.  Taking $m=Dk$
and $c=\gamma/k$ proves the first assertion.
\end{proof}

\begin{theorem}[Attractive empirical concentration]
\label{ac:concentration}
Suppose $0\le\gamma<1$, and suppose $Z_{k,\gamma}<\infty$ for the
values of $k$ under consideration.  For $f\in C^\infty(X,\mathbb R)$
and $|t|C_f<1-\gamma$,
\begin{equation}
 \log\mathbb E_{P_{k,\gamma}}
       \exp\!\left(t\sum_{i=1}^N[f(x_i)-\sigma(f)]\right)
 \le\frac{NA_ft^2}{2(1-\gamma-|t|C_f)}.
 \label{ac:mgf}
\end{equation}
In particular, if $A_f>0$, then for every $s>0$,
\begin{equation}
 P_{k,\gamma}\!\left(
    \left|\frac1N\sum_i f(x_i)-\sigma(f)\right|\ge s\right)
 \le 2\exp\!\left(
    -\frac{N(1-\gamma)s^2}{2(A_f+C_fs)}\right).
 \label{ac:tail}
\end{equation}
If $A_f=0$, then $f$ is constant and the deviation probability is
zero.  These conclusions are unconditional for $0\le\gamma<1/D$.
\end{theorem}

\begin{proof}
Let $S_f=\sum_i f(x_i)$ and
$\psi(t)=\log\mathbb E_{P_{k,\gamma}}e^{tS_f}$.
The random variable $S_f$ is bounded for fixed $N$, so differentiation
under the integral gives
$\psi''(t)=\operatorname{Var}_{P_{k,\gamma,t,f}}(S_f)$.
By Proposition~\ref{ac:variance}, $\psi'(0)=N\sigma(f)$, and
\eqref{ac:tilted-variance} bounds the second derivative on the segment
between $0$ and $t$ by $NA_f/(1-\gamma-|t|C_f)$.
Integrating twice proves \eqref{ac:mgf}.

For the upper tail in \eqref{ac:tail}, apply Chernoff's inequality
with
$t=(1-\gamma)s/(A_f+C_fs)$, which lies in the allowed interval.
Substitution in \eqref{ac:mgf} yields the exponent in
\eqref{ac:tail}.  Apply the same argument to $-f$ for the lower tail
and add the two bounds.  Lemma~\ref{ac:integrability} supplies the
unconditional coupling range.
\end{proof}

\begin{corollary}[Moment tubes and propagation of chaos]
\label{ac:tubes}
Fix real smooth functions $f_1,\ldots,f_r$ and set
$F=(f_1,\ldots,f_r)$,
$A=\max_j A_{f_j}$, and $C=\max_j C_{f_j}$.
Under the assumptions of Theorem~\ref{ac:concentration},
\begin{equation}
 P_{k,\gamma}\!\left(
 \left\|\frac1N\sum_iF(x_i)-\sigma(F)\right\|_2\ge s\right)
 \le 2r\exp\!\left(
 -\frac{N(1-\gamma)s^2}{2(rA+\sqrt r\,Cs)}\right)
 \label{ac:vector-tail}
\end{equation}
whenever $A>0$.  Thus every fixed-order projective moment tube of
width $N^{-a}$, $0<a<1/2$, has probability at least
$1-O(e^{-cN^{1-2a}})$, with constants independent of $k$ for fixed
$q,\gamma$, and moment order.

As $k\to\infty$, the empirical measure converges in probability,
in the weak topology, to $\sigma$.  For every fixed $\ell$, the
$\ell$-particle marginal converges weakly to $\sigma^{\otimes\ell}$.
\end{corollary}

\begin{proof}
If the norm in \eqref{ac:vector-tail} is at least $s$, at least one
coordinate has absolute value at least $s/\sqrt r$.  The union bound
and \eqref{ac:tail} give the displayed estimate.
For projective moments use the real coordinates
$f_j(z)=\operatorname{tr}(H_jP_z^{\otimes t})$ in a Hilbert--Schmidt
orthonormal basis of Hermitian operators on
$\operatorname{Sym}^t\mathbb C^D$, omitting the fixed trace coordinate.
These are a fixed finite family of smooth functions, and their
Euclidean norm is the Hilbert--Schmidt moment error.  Substitute
$s=N^{-a}$.

Smooth functions are uniformly dense in the continuous functions on
$X$, so the scalar bounds imply weak convergence of the empirical
measure in probability.  The measures $P_{k,\gamma}$ are permutation
invariant.  The difference between an expectation of a fixed product
$\prod_{j=1}^{\ell}f_j(x_j)$ and the corresponding expectation under
the $\ell$th power of the empirical measure is $O(\ell^2/N)$ for
bounded $f_j$: it is precisely the error from sampling indices with
replacement rather than without replacement.  Empirical convergence
and boundedness then give the product limit.  Density of finite sums
of products in $C(X^\ell)$ proves weak convergence of the marginals.
\end{proof}

\paragraph{Consequences for constrained partitions.}
Let $\Omega_N$ be one of the preceding shrinking moment tubes and
let $Q_N=\sigma^{\otimes N}(\,\cdot\mid\Omega_N)$.
Its attractive partition function satisfies the exact relation
\[
 \int e^{N\gamma E_k}\,dQ_N
 =Z_{k,\gamma}
    \frac{P_{k,\gamma}(\Omega_N)}
         {\sigma^{\otimes N}(\Omega_N)},
 \qquad
 E_k=-\frac1{kN}\log\|\Delta_k\|^2.
\]
Both probabilities tend to one at a stretched-exponential rate, by
Corollary~\ref{ac:tubes} for $\gamma$ and for $\gamma=0$.
Consequently conditioning on these tubes does not change any existing
speed-$N$ pressure limit in this coupling range.

\paragraph{Scope of the curvature argument.}
The curvature argument supplies finite-particle variance estimates,
exponential concentration of empirical observables, and propagation
of chaos for the actual attractive determinant law.  The logarithmic
energy control needed for pressure is provided by the compact-conditional
entropy method in Theorem~\ref{att:gibbs-energy-limit}.
Theorems~\ref{att:full-laplace} and \ref{att:design-laplace} establish
the full attractive Laplace principle on $\mathbb P^2$ for
$0\le\gamma<1/3$, and Corollary~\ref{att:polarized-extension}
extends the unrestricted statement to positive polarized manifolds
below their scalar section-integrability threshold.  The concentration
constants above degenerate at $\gamma=1$; extending the logarithmic
energy method toward that coupling is a separate question.
